% Figure insertions for Examples section
% Add these at appropriate locations in jss-article.tex

% Figure 2: After validation section (around line 287)
\begin{figure}[ht]
\centering
\includegraphics[width=\textwidth]{figure_2}
\caption{Validation and method comparison results. \textbf{(A)} Ground truth validation showing excellent concordance between theoretical and empirical power estimates (MAE < 0.02, R² = 0.98). \textbf{(B)} Statistical test comparison for DDA data (N=3, 2-fold effect) demonstrating superior performance of Bayes factor tests over traditional approaches. \textbf{(C)} Edge case robustness testing showing graceful degradation under challenging conditions including high missingness rates.}
\label{fig:validation}
\end{figure}

% Figure 3: After DDA case study section (around line 304)
\begin{figure}[ht]
\centering
\includegraphics[width=\textwidth]{figure_3}
\caption{DDA phosphoproteomics case study results (Arabidopsis time-course, n=2,228 peptides). \textbf{(A)} Distribution fitting identified gamma and lognormal as dominant models, with substantial heterogeneity across peptides. \textbf{(B)} Power heatmap reveals the experimental design space, with the transition from under-powered (<50\%) to well-powered (>80\%) regions clearly delineated. \textbf{(C)} Sample size requirements for 80\% of peptides to achieve 80\% power: Wilcoxon tests require N=12 per group while Bayes factor tests achieve the same threshold at N=8.}
\label{fig:dda-results}
\end{figure}

% Figure 4: After PRM analysis section (around line 320)
\begin{figure}[ht]
\centering
\includegraphics[width=\textwidth]{figure_4}
\caption{PRM targeted proteomics analysis with missing data handling (Magnaporthe oryzae, n=285 peptides). \textbf{(A)} MNAR pattern detection showing strong negative correlation (ρ = -0.42, p < 0.001) between abundance and missingness rate, indicating systematic under-sampling of low-abundance peptides. \textbf{(B)} Realistic power estimates accounting for missingness show 17\% reduction compared to complete-data scenarios. \textbf{(C)} FDR-aware analysis with Benjamini-Hochberg correction reduces median power from 45\% to 31\%, reflecting the cost of multiple testing correction in targeted panels.}
\label{fig:prm-results}
\end{figure}